\IfFileExists{stacks-project.cls}{%
\documentclass{stacks-project}
}{%
\documentclass{amsart}
}
\usepackage{amssymb}

% For dealing with references we use the comment environment
\usepackage{verbatim}
\usepackage{fancyvrb}
\newenvironment{reference}{\comment}{\endcomment}
%\newenvironment{reference}{}{}
\newenvironment{slogan}{\comment}{\endcomment}
\newenvironment{history}{\comment}{\endcomment}

% For source figures retained as assets
\usepackage{graphicx}

% For commutative diagrams we use Xy-pic
\usepackage[all]{xy}

% We use 2cell for 2-commutative diagrams.
\xyoption{2cell}
\UseAllTwocells

% We use multicol for the list of chapters between chapters
\usepackage{multicol}

% This is generally recommended for better output
\usepackage{lmodern}
\usepackage[T1]{fontenc}
\usepackage{textalpha}

% For cross-file-references
% We prefer the snapshot tmp/xref/#2.aux,
% which latexmkrc replaces atomically, so
% that documents compiled simultaneously
% never read a half-written .aux file.
% The last argument keeps links to #2.pdf.
\usepackage{xr-hyper}
\newcommand{\maybeexternaldocument}[2]{%
  \IfFileExists{#2.aux}{%
    \externaldocument[#1]{#2}%
  }{%
    \PackageWarning{stack}{Missing #2.aux; external references with prefix #1 unavailable}%
  }%
}
\let\externaldocumentorig\externaldocument
\renewcommand{\externaldocument}[2][]{%
  \IfFileExists{tmp/xref/#2.aux}{%
    \externaldocumentorig[#1]{tmp/xref/#2}[#2.pdf]%
  }{%
  \IfFileExists{#2.aux}{%
    \externaldocumentorig[#1]{#2}%
  }{%
    \PackageWarning{stack}{Missing #2.aux; external references with prefix #1 unavailable}%
  }%
  }%
}

% Package for hypertext links:
\usepackage{hyperref}

% For any local file, say "hello.tex" you want to link to please
% use \externaldocument[hello-]{hello}
\externaldocument[introduction-]{introduction}
\externaldocument[conventions-]{conventions}
\externaldocument[sets-]{sets}
\externaldocument[categories-]{categories}
\externaldocument[topology-]{topology}
\externaldocument[sheaves-]{sheaves}
\externaldocument[sites-]{sites}
\externaldocument[stacks-]{stacks}
\externaldocument[fields-]{fields}
\externaldocument[algebra-]{algebra}
\externaldocument[brauer-]{brauer}
\externaldocument[homology-]{homology}
\externaldocument[derived-]{derived}
\externaldocument[simplicial-]{simplicial}
\externaldocument[more-algebra-]{more-algebra}
\externaldocument[smoothing-]{smoothing}
\externaldocument[modules-]{modules}
\externaldocument[sites-modules-]{sites-modules}
\externaldocument[injectives-]{injectives}
\externaldocument[cohomology-]{cohomology}
\externaldocument[sites-cohomology-]{sites-cohomology}
\externaldocument[dga-]{dga}
\externaldocument[dpa-]{dpa}
\externaldocument[sdga-]{sdga}
\externaldocument[hypercovering-]{hypercovering}
\externaldocument[schemes-]{schemes}
\externaldocument[constructions-]{constructions}
\externaldocument[properties-]{properties}
\externaldocument[morphisms-]{morphisms}
\externaldocument[coherent-]{coherent}
\externaldocument[divisors-]{divisors}
\externaldocument[limits-]{limits}
\externaldocument[varieties-]{varieties}
\externaldocument[topologies-]{topologies}
\externaldocument[descent-]{descent}
\externaldocument[perfect-]{perfect}
\externaldocument[more-morphisms-]{more-morphisms}
\externaldocument[flat-]{flat}
\externaldocument[groupoids-]{groupoids}
\externaldocument[more-groupoids-]{more-groupoids}
\externaldocument[etale-]{etale}
\externaldocument[chow-]{chow}
\externaldocument[intersection-]{intersection}
\externaldocument[pic-]{pic}
\externaldocument[weil-]{weil}
\externaldocument[adequate-]{adequate}
\externaldocument[dualizing-]{dualizing}
\externaldocument[duality-]{duality}
\externaldocument[discriminant-]{discriminant}
\externaldocument[derham-]{derham}
\externaldocument[local-cohomology-]{local-cohomology}
\externaldocument[algebraization-]{algebraization}
\externaldocument[curves-]{curves}
\externaldocument[resolve-]{resolve}
\externaldocument[models-]{models}
\externaldocument[functors-]{functors}
\externaldocument[equiv-]{equiv}
\externaldocument[pione-]{pione}
\externaldocument[etale-cohomology-]{etale-cohomology}
\externaldocument[proetale-]{proetale}
\externaldocument[relative-cycles-]{relative-cycles}
\externaldocument[more-etale-]{more-etale}
\externaldocument[trace-]{trace}
\externaldocument[crystalline-]{crystalline}
\externaldocument[spaces-]{spaces}
\externaldocument[spaces-properties-]{spaces-properties}
\externaldocument[spaces-morphisms-]{spaces-morphisms}
\externaldocument[decent-spaces-]{decent-spaces}
\externaldocument[spaces-cohomology-]{spaces-cohomology}
\externaldocument[spaces-limits-]{spaces-limits}
\externaldocument[spaces-divisors-]{spaces-divisors}
\externaldocument[spaces-over-fields-]{spaces-over-fields}
\externaldocument[spaces-topologies-]{spaces-topologies}
\externaldocument[spaces-descent-]{spaces-descent}
\externaldocument[spaces-perfect-]{spaces-perfect}
\externaldocument[spaces-more-morphisms-]{spaces-more-morphisms}
\externaldocument[spaces-flat-]{spaces-flat}
\externaldocument[spaces-groupoids-]{spaces-groupoids}
\externaldocument[spaces-more-groupoids-]{spaces-more-groupoids}
\externaldocument[bootstrap-]{bootstrap}
\externaldocument[spaces-pushouts-]{spaces-pushouts}
\externaldocument[spaces-chow-]{spaces-chow}
\externaldocument[groupoids-quotients-]{groupoids-quotients}
\externaldocument[spaces-more-cohomology-]{spaces-more-cohomology}
\externaldocument[spaces-simplicial-]{spaces-simplicial}
\externaldocument[spaces-duality-]{spaces-duality}
\externaldocument[formal-spaces-]{formal-spaces}
\externaldocument[restricted-]{restricted}
\externaldocument[spaces-resolve-]{spaces-resolve}
\externaldocument[formal-defos-]{formal-defos}
\externaldocument[defos-]{defos}
\externaldocument[cotangent-]{cotangent}
\externaldocument[examples-defos-]{examples-defos}
\externaldocument[algebraic-]{algebraic}
\externaldocument[examples-stacks-]{examples-stacks}
\externaldocument[stacks-sheaves-]{stacks-sheaves}
\externaldocument[criteria-]{criteria}
\externaldocument[artin-]{artin}
\externaldocument[quot-]{quot}
\externaldocument[stacks-properties-]{stacks-properties}
\externaldocument[stacks-morphisms-]{stacks-morphisms}
\externaldocument[stacks-limits-]{stacks-limits}
\externaldocument[stacks-cohomology-]{stacks-cohomology}
\externaldocument[stacks-perfect-]{stacks-perfect}
\externaldocument[stacks-introduction-]{stacks-introduction}
\externaldocument[stacks-more-morphisms-]{stacks-more-morphisms}
\externaldocument[stacks-geometry-]{stacks-geometry}
\externaldocument[moduli-]{moduli}
\externaldocument[moduli-curves-]{moduli-curves}
\externaldocument[examples-]{examples}
\externaldocument[exercises-]{exercises}
\externaldocument[guide-]{guide}
\externaldocument[desirables-]{desirables}
\externaldocument[coding-]{coding}
\externaldocument[obsolete-]{obsolete}
\externaldocument[fdl-]{fdl}
\externaldocument[index-]{index}

% Heap Project documents
\externaldocument[k3-]{k3}
\externaldocument[torelli-]{torelli}
\externaldocument[linear-algebra-]{linear-algebra}
\externaldocument[groups-]{groups}
\externaldocument[complex-analysis-]{complex-analysis}
\externaldocument[functional-analysis-]{functional-analysis}
\externaldocument[categories-]{categories}
\externaldocument[commutative-algebra-]{commutative-algebra}
\externaldocument[ringed-spaces-]{ringed-spaces}
\externaldocument[homological-algebra-]{homological-algebra}
\externaldocument[lie-algebras-]{lie-algebras}
\externaldocument[representation-theory-]{representation-theory}
\externaldocument[smooth-manifolds-]{smooth-manifolds}
\externaldocument[differential-topology-]{differential-topology}
\externaldocument[algebraic-topology-]{algebraic-topology}
\externaldocument[differential-geometry-]{differential-geometry}
\externaldocument[algebraic-geometry-]{algebraic-geometry}
\externaldocument[symplectic-geometry-]{symplectic-geometry}
\externaldocument[complex-geometry-]{complex-geometry}
\externaldocument[hodge-theory-]{hodge-theory}
\externaldocument[%
variations-of-hodge-structures-]{%
variations-of-hodge-structures}
\externaldocument[%
homological-mirror-symmetry-]{%
homological-mirror-symmetry}
\maybeexternaldocument{deformations-}{deformations}
\externaldocument[vertex-algebras-]{vertex-algebras}
\externaldocument[springer-]{springer}
\externaldocument[infty-categories-]{infty-categories}
\externaldocument[seminars-differential-geometry-]{%
seminars-differential-geometry}
\externaldocument[seminars-complex-geometry-]{%
seminars-complex-geometry}
\externaldocument[seminars-algebraic-geometry-]{%
seminars-algebraic-geometry}
\externaldocument[seminars-others-]{%
seminars-others}
\externaldocument[%
geometric-prequantization-]{%
geometric-prequantization}
\externaldocument[surfaces-]{surfaces}
\externaldocument[birational-maps-commutative-algebra-]{birational-maps-commutative-algebra}
\externaldocument[cremona-transformations-]{cremona-transformations}
\externaldocument[derived-categories-of-sheaves-]{derived-categories-of-sheaves}
\externaldocument[geometric-invariant-theory-]{geometric-invariant-theory}
\externaldocument[geometric-stability-conditions-and-group-actions-]{geometric-stability-conditions-and-group-actions}
\externaldocument[moduli-spaces-of-sheaves-]{moduli-spaces-of-sheaves}
\externaldocument[bridgeland-stability-general-theory-]{bridgeland-stability-general-theory}
\externaldocument[hilbert-schemes-]{hilbert-schemes}
\externaldocument[noncommutative-abelian-surfaces-and-kummer-type-hyperkahler-varieties-]{noncommutative-abelian-surfaces-and-kummer-type-hyperkahler-varieties}
\externaldocument[mumford-tate-groups-in-hodge-theory-]{mumford-tate-groups-in-hodge-theory}
\externaldocument[hodge-birational-atoms-2026-]{hodge-birational-atoms-2026}
\externaldocument[xxii-egd-2026-]{%
xxii-egd-2026}
\externaldocument[pde-]{pde}
\externaldocument[probability-]{probability}

\externaldocument[linguistics-]{linguistics}
\externaldocument[genealogia-familiar-]{%
genealogia-familiar}

\externaldocument[%
16th-alga-meeting-2026-]{%
16th-alga-meeting-2026}

\externaldocument[%
iv-jornadas-lefschetz-2026-]{%
iv-jornadas-lefschetz-2026}

\externaldocument[reading-the-masters-]{%
reading-the-masters}

\externaldocument[real-analysis-]{%
real-analysis}

\externaldocument[optimal-transport-]{%
optimal-transport}

\externaldocument[history-]{history}

\externaldocument[seiberg-witten-]{seiberg-witten}

\externaldocument[history-philosophy-biology-]{history-philosophy-biology}

\externaldocument[chemistry-]{chemistry}

\externaldocument[molecular-cell-biology-]{molecular-cell-biology}

\externaldocument[microbiology-]{microbiology}

\externaldocument[protistology-]{protistology}

\externaldocument[mycology-]{mycology}

\externaldocument[botany-]{botany}

\externaldocument[zoology-]{zoology}

\externaldocument[animal-histology-]{animal-histology}

\externaldocument[animal-physiology-]{animal-physiology}

\externaldocument[systematics-biogeography-]{systematics-biogeography}

\externaldocument[ecology-natural-resources-]{ecology-natural-resources}

\externaldocument[earth-sciences-]{earth-sciences}

\externaldocument[biostatistics-]{biostatistics}

\externaldocument[mathematical-systems-biology-]{mathematical-systems-biology}

\externaldocument[cimpa-el-salvador-2024-]{cimpa-el-salvador-2024}

% Theorem environments.
%
\theoremstyle{plain}
\newtheorem{theorem}[subsection]{Theorem}
\newtheorem{proposition}[subsection]{Proposition}
\newtheorem{lemma}[subsection]{Lemma}
\newtheorem{corollary}[subsection]{Corollary}

\theoremstyle{definition}
\newtheorem{definition}[subsection]{Definition}
\newtheorem{example}[subsection]{Example}
\newtheorem{exercise}[subsection]{Exercise}
\newtheorem{situation}[subsection]{Situation}

\theoremstyle{remark}
\newtheorem{remark}[subsection]{Remark}
\newtheorem{remarks}[subsection]{Remarks}

\numberwithin{equation}{subsection}

% Macros
%
\def\lim{\mathop{\mathrm{lim}}\nolimits}
\def\colim{\mathop{\mathrm{colim}}\nolimits}
\def\Spec{\mathop{\mathrm{Spec}}}
\def\Hom{\mathop{\mathrm{Hom}}\nolimits}
\def\Ext{\mathop{\mathrm{Ext}}\nolimits}
\def\SheafHom{\mathop{\mathcal{H}\!\mathit{om}}\nolimits}
\def\SheafExt{\mathop{\mathcal{E}\!\mathit{xt}}\nolimits}
\def\Sch{\mathit{Sch}}
\def\Mor{\mathop{\mathrm{Mor}}\nolimits}
\def\Ob{\mathop{\mathrm{Ob}}\nolimits}
\def\Sh{\mathop{\mathit{Sh}}\nolimits}
\def\NL{\mathop{N\!L}\nolimits}
\def\CH{\mathop{\mathrm{CH}}\nolimits}
\def\proetale{{pro\text{-}\acute{e}tale}}
\def\etale{{\acute{e}tale}}
\def\QCoh{\mathit{QCoh}}
\def\Ker{\mathop{\mathrm{Ker}}}
\def\Im{\mathop{\mathrm{Im}}}
\def\Coker{\mathop{\mathrm{Coker}}}
\def\Coim{\mathop{\mathrm{Coim}}}

% Boxtimes
%
\DeclareMathSymbol{\boxtimes}{\mathbin}{AMSa}{"02}

%
% Macros for moduli stacks/spaces
%
\def\QCohstack{\mathcal{QC}\!\mathit{oh}}
\def\Cohstack{\mathcal{C}\!\mathit{oh}}
\def\Spacesstack{\mathcal{S}\!\mathit{paces}}
\def\Quotfunctor{\mathrm{Quot}}
\def\Hilbfunctor{\mathrm{Hilb}}
\def\Curvesstack{\mathcal{C}\!\mathit{urves}}
\def\Polarizedstack{\mathcal{P}\!\mathit{olarized}}
\def\Complexesstack{\mathcal{C}\!\mathit{omplexes}}
% \Pic is the operator that assigns to X its picard group, usage \Pic(X)
% \Picardstack_{X/B} denotes the Picard stack of X over B
% \Picardfunctor_{X/B} denotes the Picard functor of X over B
\def\Pic{\mathop{\mathrm{Pic}}\nolimits}
\def\Picardstack{\mathcal{P}\!\mathit{ic}}
\def\Picardfunctor{\mathrm{Pic}}
\def\Deformationcategory{\mathcal{D}\!\mathit{ef}}


\begin{document}

\title{Systematics and Biogeography}
\maketitle

\phantomsection
\label{section-phantom}

\tableofcontents

\section{What this chapter covers}
\label{section-overview}

Systematics: the schools of classification,
characters and homology, parsimony and the
measures of fit and support of a
cladogram, phenetic clustering, ranks and
nomenclature, species concepts and modes
of speciation. Biogeography: its history,
Darwin's dispersal experiments, the
schools of historical biogeography, island
biogeography, quantitative methods, the
molecular clock, and the biogeographic
provinces of Mexico.

\textbf{Sources of these notes.}
The first half rests on the course
\emph{Sistem\'atica} (second semester,
Facultad de Ciencias, UNAM, spring 2014),
the second on \emph{Biogeograf\'ia} (third
semester, autumn 2014). What survives is
Daniel's work: exam outlines for both
courses; lab exercises (coding characters
of animals, fishes and books; a skeletal
matrix of five ornithopod dinosaurs; a
phenogram computed by hand); a team essay
on the species concept; slides on the
reptiles of Mexico; summaries of two
papers on Mexican cloud forests; an
opinion on chapters XI and XII of the
\emph{Origin}; a short essay on the
molecular clock; a comparison of maps of
Mexican biotic provinces; and a team
project on the spread of the monk
parakeet. The facts were checked and
completed with Simpson's \emph{Plant
Systematics} \cite{simpson2006plant},
Hickman, Roberts and Larson's \emph{Animal
Diversity} \cite{hickman2002animal},
Singh's \emph{Plant Systematics}
\cite{singh2010plant}, Mayr's \emph{What
Evolution Is} \cite{mayr2002evolution},
Darwin's \emph{Origin}
\cite{darwin2008origin}, Hubbell's
\emph{Unified Neutral Theory}
\cite{hubbell2001unified} and Morrone's
\emph{Biogeograf\'ia de Am\'erica Latina y
el Caribe} \cite{morrone2001biogeografia}.

\section{Schools of systematics}
\label{section-schools}

The course's history ran from folk
taxonomy and Greek essentialism (the
archetype) through Linnaeus and Darwin to
three modern schools.

\textbf{Evolutionary systematics}
(Simpson, Mayr) uses two criteria: common
descent and the amount of adaptive change.
A branch that enters a new adaptive zone is
a \emph{grade} and earns a taxon; the
broader the zone, the higher the rank.
Penguins are a grade within birds.
Evolutionary taxa may be monophyletic or
paraphyletic: the traditional family
Pongidae (orangutan, gorilla, chimpanzee)
leaves out humans, although gorillas and
chimpanzees are closer to humans than to
orangutans
\cite[Chapter 4]{hickman2002animal}.

\textbf{Phenetics} (numerical taxonomy)
groups taxa by overall similarity, without
regard to phylogeny. It may group taxa by
shared ancestral states, so its
phenograms can differ from cladograms
\cite[Chapters 1 and 2]{simpson2006plant}.
It had little impact on animal
classification
\cite[Chapter 4]{hickman2002animal}.

\textbf{Phylogenetic systematics}
(cladistics), founded by Willi Hennig
(main ideas published shortly before 1950,
English translation 1966), takes the
genealogical branching pattern as the only
basis of classification, because it alone
gives an unambiguous hierarchy, and uses
shared derived characters
(synapomorphies) as the only evidence for
it. Only monophyletic groups are named
\cite[Chapter 2]{simpson2006plant}.

\section{Characters, homology and groups}
\label{section-characters}

\begin{definition}
\label{definition-character}
A \emph{character} is a feature that
varies among the taxa studied; its
variants are \emph{character states}.
Characters may be morphological,
chemical, molecular, behavioural or
biogeographic.
\end{definition}

\begin{definition}
\label{definition-homology}
Two features are \emph{homologous} if
their similarity is due to common
ancestry. Similarity not due to common
ancestry is \emph{homoplasy}: by
\emph{convergence} (different ancestral
conditions), \emph{parallelism} (the same
ancestral condition) or \emph{reversal} (a
return to the ancestral state).
\end{definition}

Homology is a hypothesis, tested by the
cladistic analysis itself: a feature that
turns out to have arisen several times was
not homologous
\cite[Chapter 2]{simpson2006plant}.

\begin{definition}
\label{definition-apomorphy}
Relative to a monophyletic group, an
ancestral state is a \emph{plesiomorphy}
and a derived state an \emph{apomorphy}. A
\emph{synapomorphy} is an apomorphy shared
by two or more taxa, an
\emph{autapomorphy} one found in a single
taxon, and a \emph{symplesiomorphy} a
shared plesiomorphy.
\end{definition}

The terms describe states relative to a
group, not taxa: a state can be derived
within a large group and ancestral within
one of its subgroups
\cite[Chapter 2]{simpson2006plant}.

\begin{definition}
\label{definition-monophyly}
A group is \emph{monophyletic} if it
contains a common ancestor and all its
descendants, \emph{paraphyletic} if it
contains the ancestor and some but not all
descendants, and \emph{polyphyletic} if it
does not contain the most recent common
ancestor of its members.
\end{definition}

Monophyletic groups are diagnosed by
synapomorphies, paraphyletic ones by
symplesiomorphies, polyphyletic ones by
convergences. Birds and mammals together
(``Homeothermia'') would be polyphyletic
\cite[Chapter 4]{hickman2002animal}.

\textbf{Coding.} States are coded by
nonnegative integers, the outgroup state
usually 0. Binary characters take 0
(plesiomorphic) or 1 (apomorphic).
Multistate characters are \emph{ordered}
(additive: from 0 to 2 one must pass
through 1) or \emph{unordered}; the usual
practice is unordered unless there is
strong evidence for an order. A
\emph{step matrix} gives the cost of each
change; for DNA, a transition (purine to
purine, pyrimidine to pyrimidine) may cost
1 step and a transversion (purine to
pyrimidine) 5
\cite[Chapter 2]{simpson2006plant}.
Weighting can be a priori (by judged
complexity, criticized as subjective),
by compatibility (only characters without
homoplasy), or \emph{successive}: build a
tree with equal weights, reweight each
character by its fit (often by the
rescaled consistency index), rebuild.

\textbf{Polarity: the outgroup rule.} The
outgroup is a group close to, but outside,
the study group (ingroup).
\begin{itemize}
\item A state found both in the ingroup and
in the outgroup is plesiomorphic.
\item A state found only in the ingroup is
apomorphic.
\item If the outgroup has all the states,
use a more distant outgroup.
\end{itemize}
The rationale is parsimony: the other
choice needs more changes. Teeth are
present in amphibians and bony fishes, so
their absence in birds is derived within
amniotes \cite[Chapter 4]{hickman2002animal}.
The course also listed weaker criteria:
fossils show the plesiomorphic state,
``common is primitive'', and functional
organs are plesiomorphic relative to
vestigial ones.

\section{Parsimony and tree search}
\label{section-parsimony}

Parsimony prefers the cladogram that
needs the fewest character state changes
(steps). The cost rules differ:
\begin{center}
\begin{tabular}{ll}
Criterion & Changes allowed \\
\hline
Wagner & ordered; $0 \rightarrow 2$
costs 2 \\
Fitch & unordered; every change 1 \\
Dollo & gain once, losses free \\
Camin--Sokal & no reversals \\
generalized & any step matrix
\end{tabular}
\end{center}
Dollo's rule reflects the idea that a
complex structure evolves once but can be
lost many times; Camin--Sokal is the
irreversible step matrix of
\cite[Chapter 2]{simpson2006plant}.

Methods of search, as listed in the
course:
\begin{itemize}
\item \emph{Hennigian argumentation}:
group only by synapomorphies, combine
compatible characters, assume homology
unless there is evidence against it.
\item \emph{Wagner method}: add taxa one
at a time, starting with the one most
similar to the outgroup, joining each new
taxon to the most similar branch or
hypothetical ancestor; practical for about
twenty taxa.
\item \emph{Exhaustive search}: build all
trees and keep the shortest.
\item \emph{Branch and bound}: grow trees
taxon by taxon and abandon any partial
tree already longer than the best complete
tree found.
\item \emph{Branch swapping} (heuristic):
rearrange a tree locally (swapping
branches around an internal branch) or
globally (cutting and regrafting
subtrees); the \emph{ratchet} handles very
large matrices.
\end{itemize}

\section{Fit, support and consensus}
\label{section-support}

Let $s$ be the number of steps on a tree,
$m$ the minimum number of steps the data
could need on any tree, and $g$ the
maximum, the steps needed on a completely
unresolved tree. Then
\begin{equation}
\label{equation-ci-ri}
\begin{aligned}
CI &= \frac{m}{s}, \\
RI &= \frac{g-s}{g-m}, \\
RC &= CI \cdot RI
\end{aligned}
\end{equation}
are the \emph{consistency index}, the
\emph{retention index} and the
\emph{rescaled consistency index}. Values
near 1 mean little homoplasy. The CI can
be computed per character, cannot reach 0,
and is inflated by autapomorphies; the RI
depends on the number of taxa
\cite[Chapter 2]{simpson2006plant}.

\textbf{Support.} The \emph{bootstrap}
resamples characters with replacement to
build matrices of the original size,
reruns the analysis many times (100 or
more) and reports how often each clade
appears; 70\% or more is usually taken as
robust. The \emph{jackknife} resamples
without replacement, giving smaller
matrices. The \emph{decay index} or
\emph{Bremer support} of a clade is the
number of extra steps needed before the
clade disappears from the shortest trees
\cite[Chapter 2]{simpson2006plant}.

\textbf{Consensus trees} summarize several
equally parsimonious trees. The
\emph{strict} consensus keeps only the
clades present in all of them, collapsing
the rest into polytomies; the \emph{50\%
majority-rule} consensus keeps clades
present in at least half
\cite[Chapter 2]{simpson2006plant}. The
course also listed semistrict, Adams and
Nelson consensus trees, and supertrees.

\textbf{Likelihood and Bayes.}
\emph{Maximum likelihood} chooses the tree
under which a chosen model of evolution
gives the observed data the highest
probability; the model can, for instance,
make transitions likelier than
transversions and let rates vary.
\emph{Bayesian inference} computes the
posterior probability of trees, sampled by
Markov chain Monte Carlo; a clade with
posterior probability 95\% or more is
considered well supported, although these
values are often high
\cite[Chapter 2]{simpson2006plant}. For a
tree $H$ and data $D$,
$$
P(H \mid D) =
\frac{P(D \mid H)\, P(H)}{P(D)},
$$
where $P(D \mid H)$ is the likelihood,
$P(H)$ the prior and $P(D)$ the marginal
likelihood, the same for all trees.

\begin{example}
\label{example-parsimony}
Take an outgroup O, three taxa A, B, C,
and four binary characters:
\begin{center}
\begin{tabular}{lcccc}
 & 1 & 2 & 3 & 4 \\
\hline
O & 0 & 0 & 0 & 0 \\
A & 1 & 1 & 0 & 0 \\
B & 1 & 1 & 1 & 0 \\
C & 0 & 0 & 1 & 1
\end{tabular}
\end{center}
Rooted at O there are three trees. On
((A,B),C) characters 1, 2 and 4 change
once and character 3 twice: length 5. On
((B,C),A) characters 1 and 2 need two
steps each: length 6. On ((A,C),B)
characters 1, 2 and 3 need two steps:
length 7. So ((A,B),C) is the most
parsimonious tree, with $s = 5$ and
$m = 4$. For $g$, each binary character
needs on an unresolved tree as many steps
as the taxa in its rarer state:
$2+2+2+1 = 7$.
Hence
$$
CI = \frac{4}{5} = 0.8, \quad
RI = \frac{7-5}{7-4} \approx 0.67, \quad
RC \approx 0.53.
$$
Character 4 is an autapomorphy of C: it
adds 1 to both $m$ and $s$ but says
nothing about grouping; without it
$CI = 3/4$. The clade (A,B) is lost on
the next best tree, one step longer, so
its Bremer support is 1.
\end{example}

\section{Phenetics}
\label{section-phenetics}

Phenetic clustering starts from a matrix
of similarities between operational
taxonomic units (OTUs). For binary data,
with $a$ the number of characters present
in both OTUs, $d$ absent in both, and $u$
the mismatches (present in only one), the
\emph{simple matching coefficient} and the
\emph{Jaccard coefficient} are
$$
S_{SM} = \frac{a+d}{a+d+u}, \qquad
S_J = \frac{a}{a+u}.
$$
Jaccard ignores shared absences, which
suits presence/absence data. In
\emph{UPGMA} (unweighted pair-group method
with arithmetic averages) the two most
similar clusters are merged and the
similarity of the new cluster to each
other cluster is the average over all
pairs of members
\cite[Chapter 8]{singh2010plant}.

\begin{example}
\label{example-upgma}
In the course's lab, Jaccard similarities
among four hypothetical OTUs W, X, Y, Z
were:
\begin{center}
\begin{tabular}{lccc}
 & W & X & Y \\
\hline
X & 0.20 & & \\
Y & 0.17 & 0.80 & \\
Z & 0.33 & 0.00 & 0.17
\end{tabular}
\end{center}
UPGMA joins X and Y at 0.80. The new
similarities are $(0.20+0.17)/2 \approx
0.18$ for W and $(0.00+0.17)/2 \approx
0.08$ for Z, so W and Z join next at
0.33. The two clusters join at the mean of
the four cross pairs,
$(0.20+0.17+0.00+0.17)/4 \approx 0.13$.
The phenogram is ((X,Y),(W,Z)).
\end{example}

\section{Ranks and nomenclature}
\label{section-nomenclature}

\textbf{Ranks.} The seven mandatory
ranks of the animal kingdom are kingdom,
phylum, class, order, family, genus and
species; intermediate ranks (subphylum,
subclass, suborder, subfamily) are added
as needed. A species name is a binomial,
genus plus specific epithet, in italics; a
subspecies takes a trinomial, as in
\emph{Turdus migratorius achrustera}. The
author and year that follow a name
(\emph{Didelphis marsupialis} Linnaeus,
1758) are not part of it; parentheses
mark that the species has since moved to
another genus. Formal subspecies have
lost favour, since their boundaries are
rarely sharp
\cite[Chapter 4]{hickman2002animal}.

\textbf{Types.} A name rests on a type.
The \emph{holotype} is the single specimen
designated in the original publication; an
\emph{isotype} is a duplicate of it (in
botany); a \emph{lectotype} is chosen
later from the original material when no
holotype was designated or it is lost; a
\emph{neotype} is chosen from new material
when all original material is lost; a
\emph{syntype} is any specimen cited when
no holotype was designated; a
\emph{paratype} is a specimen cited that
is none of these. A genus name is tied to
a type species, a family name to a type
genus \cite[Chapter 16]{simpson2006plant}.

\textbf{Rules.} \emph{Priority}: of
competing names the earliest published is
correct. \emph{Synonyms} are rejected
names for the same taxon, either
illegitimate or rejected by taxonomic
judgement; each taxon has one correct
name. \emph{Homonyms} are identical names
based on different types; the later
homonym is illegitimate. A
\emph{basionym} is the original name that
supplies the epithet when a taxon is moved
or changes rank
\cite[Chapter 16]{simpson2006plant}. The
course added monotypy (the type species is
the only species of its genus) and
tautonymy (genus and epithet identical).

\textbf{Example: the reptiles of
Mexico.} The course's slides used the
ranks of the lab manual:
\begin{center}
\begin{tabular}{ll}
Rank & Taxa \\
\hline
class & Reptilia \\
superorder & Lepidosauria, Archosauria \\
order & Testudines, Squamata, \\
 & Crocodylia \\
suborder & Pleurodira, Cryptodira \\
 & (turtles); Lacertilia, \\
 & Amphisbaenia, Serpentes
\end{tabular}
\end{center}
Flores-Villela and Garc\'ia-V\'azquez
(\emph{Revista Mexicana de
Biodiversidad}, 2014), read in the course,
count for Mexico (to October 2013) 864
species in 159 genera and 40 families,
8.7\% of the world's reptiles: 417
lizards, 393 snakes, 48 turtles, 3
amphisbaenians and 3 crocodilians; 493
species (57\%) are endemic. The richest
states are Oaxaca (262), Chiapas (220) and
Veracruz (200). The \emph{Systema
Naturae} already lists several Mexican
reptiles, 13 of them still valid. The
paper stresses that reptiles are not a
natural group (crocodiles share
synapomorphies with birds), that the major
groups are better left without Linnaean
ranks, and that lizards are paraphyletic,
since snakes and amphisbaenians are nested
among them. Reptiles as traditionally
recognized are paraphyletic
\cite[Chapter 4]{hickman2002animal}.

\section{Species and speciation}
\label{section-species}

The course listed many species concepts:
typological, phenetic, nominalist,
morphological, biological, evolutionary,
the cohesion or selection concept, and
others. The reading by Llorente and
Mich\'an frames two old debates behind
them: realism against nominalism, and
fixism against evolutionism.
\begin{itemize}
\item \emph{Typological}: a species is
defined by fixed essential features,
usually morphological; discarded after
Darwin.
\item \emph{Biological} (Dobzhansky,
Mayr): ``a reproductive community of
populations (reproductively isolated from
others) that occupies a specific niche in
nature'' (Mayr, 1983). It does not apply
to asexual organisms.
\item \emph{Evolutionary} (Simpson): a
single lineage of ancestor--descendant
populations with its own identity,
evolutionary tendencies and historical
fate; it covers asexual forms.
\item \emph{Phylogenetic}: the smallest
diagnosably distinct grouping with a
parental pattern of ancestry and descent;
it yields the most species.
\end{itemize}
\cite[Chapter 4]{hickman2002animal}.
Mayr accepts only two concepts, the
typological and the biological, and reads
the others as criteria for delimiting
species taxa. Sibling species, such as the
six species once lumped as the malaria
mosquito \emph{Anopheles maculipennis},
defeat the typological concept
\cite[Chapter 8]{mayr2002evolution}.

\textbf{Speciation.} Anagenesis is change
within a lineage, cladogenesis the
splitting of lineages.
\begin{itemize}
\item \emph{Allopatric}: a population
acquires isolating mechanisms while
geographically isolated. In
\emph{dichopatric} speciation a new
barrier splits the range (vicariance); in
\emph{peripatric} speciation a small
founder population beyond the edge of the
range diverges
\cite[Chapter 9]{mayr2002evolution}.
\item \emph{Sympatric}: without
geographic isolation, as when insects
shift to a new host plant and mate on it,
or in cichlid fishes of small crater lakes
in Cameroon
\cite[Chapter 9]{mayr2002evolution}.
\item The course also listed
\emph{parapatric} speciation (adjacent
populations, with loss of the
intermediates of a cline) and
\emph{stasipatric} speciation (by
chromosomal change).
\end{itemize}

\textbf{The team essay.} Daniel's team
argued, with Leibniz's principles of
identity and continuity, that each
organism is unique while forms grade into
one another, so that every species concept
cuts a continuous process arbitrarily. Its
examples were a DNA barcode library of the
butterflies of Peninsular Malaysia, where
barcodes and morphology disagreed on
species and subspecies, and a paper
defending subspecies only as evolutionarily
significant units. It floated the idea of a
species as an emergent property of a
complex system, and in the end preferred,
for practical work such as conservation,
the \emph{pragmatic} concept quoted by
Llorente and Mich\'an: the smallest group
consistently and persistently distinct from
others and distinguishable by ordinary
means.

\section{History of biogeography}
\label{section-biogeography-history}

Biogeography studies the distribution of
living things in space and time and its
causes; distribution depends on both
ecology and history. Historical
biogeography works at geological time
scales and seeks relationships between
areas; ecological biogeography works at
short times and local scales.

The course's history:
\begin{itemize}
\item a single centre of creation and
dispersal (Eden, the Flood);
\item Acosta (1590): the American biota
reached the New World by land bridges,
such as the strait of Ani\'an in the
north;
\item Linnaeus: distribution follows
ecology, as in the zones of a mountain;
\item Buffon: regions with the same
climate hold different species, so
ecology is not enough;
\item de Candolle: historical versus
ecological causes;
\item Sclater (1858): six regions from the
distribution of birds, still in use:
Palearctic, Nearctic, Neotropical,
Ethiopian, Oriental, Australian;
\item Wallace (1876): \emph{The Geographical
Distribution of Animals}, built on
Sclater's regions;
\item Darwin: dispersal over a stable
geography, against land bridges and
multiple centres of creation;
\item Croizat (\emph{Space, Time, Form},
1964): the earth and its biota evolve
together.
\end{itemize}

\textbf{Darwin's experiments}
(\emph{Origin}, chapters XI and XII). To
test ``occasional means of
distribution'', Darwin soaked seeds in sea
water: 64 of 87 kinds germinated after 28
days, and a few survived 137 days. Most
seeds sank, but dried plants with ripe
fruit floated longer: 18 of 94 floated for
more than 28 days. He concluded that the
seeds of about 14 in 100 plants could be
carried 924 miles by currents averaging 33
miles a day and still germinate. He grew
537 plants from three tablespoons of pond
mud ($6\frac{3}{4}$ ounces when dry), and
argued that wading birds carry such mud on
their feet \cite[Chapters XI and
XII]{darwin2008origin}. In Daniel's
opinion essay, Darwin assumed a stable
geography and so could not explain
distribution fully, but his dispersal was
far more reasonable than land bridges and
multiple creations, and he presented his
conclusions carefully as beliefs based on
experiment.

\section{Schools of historical biogeography}
\label{section-schools-biogeography}

\textbf{Dispersalism} (Darwin, Wallace,
Simpson). Each taxon arises in a centre of
origin and spreads. Dispersal is
ecological (within the range) or
historical (range expansion), active or
passive. Routes are \emph{corridors}
(open to many species), \emph{filters}
(passable by some) or chance routes.
Method: find the centre of origin, find
the routes, reconstruct the history.
Criticisms: no precise method, no
criterion for the centre of origin,
stable geography assumed, histories unique
to each taxon (no general patterns), abuse
of land bridges.

\textbf{Phylogenetic biogeography}
(Hennig, Brundin, Ball). It rests on
continental drift, Hennig's cladistics and
Croizat's critique of dispersalism, and
studies one monophyletic group at a time.
Two rules:
\begin{itemize}
\item \emph{deviation rule}: when a
lineage splits, one daughter differs more
and the other stays more like the
ancestor;
\item \emph{progression rule}: change in
space runs parallel to change in
characters, so the species with the most
plesiomorphic characters live in the area
occupied first, taken as the centre of
origin.
\end{itemize}
Method: build the cladogram of the group,
replace species by their areas, place the
centre of origin, compare with geological
history to decide between vicariance and
dispersal, date with fossils (minimum
ages) and barriers (maximum ages); jump
dispersal is admitted when geology and
cladogram disagree.

\textbf{Cladistic vicariance biogeography}
(Nelson, Platnick, Rosen) compares several
monophyletic groups living in areas of
endemism. Congruent patterns imply a
common history, most probably vicariance:
populations split by a new barrier.
Method: replace the terminal taxa of each
cladogram by their areas (taxon-area
cladograms), solve the problems, then
combine them into a general area cladogram
by component analysis or parsimony. The
trivial case has one endemic taxon per
area; the problems are widespread taxa,
missing areas and redundant distributions.
The school grew with plate tectonics:
Wegener proposed drift in 1915, and a
mechanism was accepted only at the end of
the 1960s.

\textbf{Panbiogeography} (Croizat) needs
no phylogeny, rejects dispersal from
centres of origin, and gives priority to
the geography of distributions: biotas
spread in a phase of mobility, then
differentiate in place when the land
changes.
\begin{itemize}
\item \emph{Individual track}: the line
joining the localities of a taxon by the
shortest total length (a minimum spanning
tree).
\item \emph{Generalized track}: several
coinciding individual tracks, a sign of
shared history.
\item \emph{Node}: where generalized tracks
meet; read as zones of tectonic
convergence and as conservation
priorities.
\item \emph{Baseline}: the barrier or
geographic feature that tracks cross, used
to orient them; a \emph{mass centre}, the
area of greatest diversity, also orients
them.
\end{itemize}
These schools are described from the
course's exams and the reading by
Llorente, Papavero and Bueno on the
history of biogeography.
Morrone uses generalized tracks to split
South America: taxa of the west connect to
Australia and New Zealand, those of the
east to the Old World tropics
\cite{morrone2001biogeografia}.

\section{Islands, quantitative methods}
\label{section-islands}

\textbf{Island biogeography.} MacArthur and
Wilson (1963, 1967) explained why islands
have fewer species than mainland areas of
the same size: smaller populations go
extinct more easily and recolonize more
slowly. The number of species $S$ on an
island settles where the immigration rate
$I$, falling as $S$ grows, meets the
extinction rate $E$, rising as $S$ grows:
$$
I(S^{*}) = E(S^{*}).
$$
Distance from the mainland lowers
immigration and small size raises
extinction; the only state variable is
the number of species, and all species are
treated as equivalent
\cite[Chapter 1]{hubbell2001unified}. The
course added four stages (initial,
competitive decline, reorganization,
evolutionary) and four weaknesses: species
can arise on the island, islands are not
homogeneous, colonization and extinction
are not independent, and history is
ignored.

\textbf{Quantitative biogeography.}
Choose a study area, divide it into
operational geographic units (grid cells),
build a presence/absence matrix (taxa by
units), compute diversity per unit and
similarity between units (S{\o}rensen's
index for presence data, Bray--Curtis for
abundances), and cluster.
\emph{Parsimony analysis of endemicity}
(PAE) treats areas as taxa and taxa as
characters, roots the tree with a
hypothetical area with all taxa absent,
and finds the most parsimonious cladogram
of areas; taxa present in a single area,
like autapomorphies, are uninformative. In
a paper summarized in the course, Luna
Vega and coauthors (\emph{Journal of
Biogeography}, 1999) ran PAE on 1267
species of vascular plants in 24 Mexican
cloud forests and found one tree of 3581
steps, with $CI = 0.354$ and
$RI = 0.291$, grouping the forests into
five clades that do not match Rzedowski's
floristic provinces.

\textbf{Molecular clock.} Zuckerkandl and
Pauling (1962) showed that many molecules
change at a roughly constant rate. Each
molecule has its own rate, calibrated
with dated fossils or vicariance events
\cite[Chapter 2]{mayr2002evolution}. If a
gene diverges at rate $r$ per lineage and
two species differ by a fraction $d$, their
common ancestor lived
$$
t = \frac{d}{2r}
$$
ago: with $r = 2\%$ per million years and
$d = 10\%$, $t = 2.5$ million years.
Rates vary between lineages (rodents
evolve faster than primates)
\cite[Chapter 10]{mayr2002evolution}.

\section{Biogeography of Mexico}
\label{section-mexico}

Morrone divides Latin America into
regions, subregions and provinces. Mexico
holds:
\begin{itemize}
\item Nearctic region: provinces of
California, Baja California, Sonora,
Altiplano Mexicano and Tamaulipas;
\item Neotropical region, Caribbean
subregion: Sierra Madre Occidental,
Sierra Madre Oriental, Eje Volc\'anico
Transmexicano, Depresi\'on del Balsas,
Sierra Madre del Sur, Costa Pac\'ifica
Mexicana, Golfo de M\'exico, Pen\'insula
de Yucat\'an and Chiapas.
\end{itemize}
Halffter's \emph{Mexican Transition Zone},
in this scheme the Mexican and Central
American part of the Caribbean subregion,
lies between the Nearctic and Neotropical
regions. Rzedowski saw its conifer and oak
forests as a southern extension of the
Nearctic, although most of its plants are
Neotropical. A panbiogeographic analysis
of Theaceae places a node in Mesoamerica,
where the Caribbean generalized track
meets an Asian--Mesoamerican one. The
Transmexican Volcanic Belt, in central
Mexico, is dominated by pine--oak forest,
with alpine grassland near the summits of
the great volcanoes
\cite{morrone2001biogeografia}.

In a course exercise Daniel compared four
CONABIO maps of Mexico (floristic,
mammal, herpetofaunal and morphotectonic
provinces). All show north--south
divisions in the north, set by the Sierras
Madre Occidental and Oriental, and
east--west divisions in the south, set by
the Transmexican Volcanic Belt, with the
Sierra Madre de Chiapas as another
barrier. A second paper summarized in the
course, a track analysis of 967 species of
vascular plants in the cloud forests of
Hidalgo by Luna Vega and coauthors
(\emph{Diversity and Distributions},
2000), found a node at Chapulhuac\'an and
proposed it as a conservation priority.

\textbf{The monk parakeet project.} The
course's final team project (November
2014) traced the spread of the monk
parakeet, \emph{Myiopsitta monachus},
native to southern Brazil, Uruguay,
Bolivia and northern Argentina and the
only parrot that builds its own nests.
It was first recorded free in New York in
1967, in Barcelona in 1975 and in England
in 1987; Chile imported about 13\,000
legally in 1995--1998. Mexico imported
about 125\,000 between 2005 and 2010 and
became the largest importer. The first
Mexican sightings were in the Valley of
Mexico in 2008, before those in Chihuahua
and Baja California (about 2011). The team
argued that the spread followed the pet
trade, not natural dispersal from the
north.

\bibliography{my}
\bibliographystyle{amsalpha}

\end{document}
